%% sec7_3.tex — Section 7.3: Discussion and Bibliography
%% Chapter 7 of "A First Course in Monte Carlo Methods"
%% Sanz-Alonso & Al-Ghattas, University of Chicago

\section{Discussion and Bibliography}
\label{sec:7.3}

The simulated annealing algorithm was introduced in \cite{124}. Geman and Geman \cite{82}
conjectured and Gidas \cite{87} rigorously showed convergence with logarithmic tempering in
finite state space. Later, Hajek \cite{103} showed that choosing the inverse temperatures via
logarithmic tempering is a sufficient and necessary condition for global convergence. Further
early results were reviewed and developed in \cite{228, 221, 1, 2}. For more recent
expositions, we refer to \cite{109, 149}.

Simulated and parallel tempering were introduced in \cite{220, 154}, and further developed in
\cite{85, 114}. For a review on parallel tempering, we refer to \cite{66}. The question of how
to choose the temperatures has been widely studied, and recent theory suggests that a geometric
sequence of temperature ratios is optimal under mild assumptions \cite{63}. Another important
question is how often one should propose swaps between chains. In a continuous time setting,
recent theory indicates the potential advantage of frequent swaps, formalized via the infinite
swapping limit for parallel tempering \cite{62}.

Augmenting the target to speed-up mixing and running multiple chains that communicate with each
other are two important ideas that underpin many sampling algorithms and convergence diagnostics
(see e.g.\ Hamiltonian Monte Carlo in Chapter~8, Box--Muller sampling method in Chapter~3,
slice sampler in Chapter~5, Gelman and Rubin's diagnostic in Chapter~4). Likewise, introducing
a sequence of auxiliary target densities to gradually reach a desired target is a powerful idea
that underpins many sequential Monte Carlo algorithms (see particle filters in Chapter~9).
Algorithms that rely on tempered transitions for sampling multi-modal distributions are
discussed in \cite{168}.
